\documentclass[11pt]{amsart}
\usepackage{bm}
\usepackage{amsmath,amsfonts, amsthm, amssymb}
\usepackage[abbrev,backrefs]{amsrefs}
%\usepackage{showlabels}
\usepackage{microtype}
\usepackage{enumerate}
\usepackage{graphicx}
%\usepackage{braket}
\graphicspath{{../../figures/}}
\usepackage[onehalfspacing]{setspace}

%%%% for widecheck command
% code from mathabx.sty and mathabx.dcl
\DeclareFontFamily{U}{mathx}{\hyphenchar\font45}
\DeclareFontShape{U}{mathx}{m}{n}{
      <5> <6> <7> <8> <9> <10>
      <10.95> <12> <14.4> <17.28> <20.74> <24.88>
      mathx10
      }{}
\DeclareSymbolFont{mathx}{U}{mathx}{m}{n}
\DeclareFontSubstitution{U}{mathx}{m}{n}
\DeclareMathAccent{\widecheck}{0}{mathx}{"71}

\setlength{\oddsidemargin}{0pt}
\setlength{\evensidemargin}{0pt}
\setlength{\textwidth}{6.0in}
\setlength{\topmargin}{0in}
\setlength{\textheight}{8.5in}

\setlength{\parindent}{0in}
\setlength{\parskip}{5px}

\usepackage{scalerel,stackengine}
\def\apeqA{\SavedStyle\sim}
\def\apeq{\setstackgap{L}{\dimexpr.5pt+1.5\LMpt}\ensurestackMath{%
  \ThisStyle{\mathrel{\Centerstack{{\apeqA} {\apeqA} {\apeqA}}}}}}

\def\bra#1{\mathinner{\langle{#1}|}}
\def\ket#1{\mathinner{|{#1}\rangle}}
\def\braket#1{\mathinner{\langle{#1}\rangle}}
\def\Braket#1{\langle{#1}\rangle}

\def\Real{{\mathbf R}}
\def\Sphere{{\mathbf S}}
\def\Schwartz{{\mathcal{S}}}
\def\Lebesgue{{\mathcal{L}}}
\def\Hausdorff{{\mathcal{H}}}
\def\Prob{{\mathbf P}}
\def\Var{{\mathbf{Var}\,}}
\def\Expec{{\mathbb E\,}}
\def\xp{{\mathbb E\,}}
\def\RWExpec{{\mathbb E\,}}
\def\eps{{\varepsilon}}
\def\One{{\mathbf{1}}}

\def\PhaseSpace{{\Real^{2d}}}
\def\Evol{{\mathcal{E}}}
\def\xop{{\mathbf{x}}}
\def\pop{{\mathbf{p}}}
\DeclareMathOperator{\Rec}{Rec}
\DeclareMathOperator{\Err}{Err}
\def\lsim{{\,\lesssim\,}}
\def\llog{{\,\lessapprox\,}}

% the Hilbert space
\def\Hilbert{{\mathcal{H}}}

\def\bbZ{{\mathbb{Z}}}
\def\bbN{{\mathbb N}}

\newcommand{\Wigner}[1]{{\mathcal{W}_{#1}}}
\newcommand{\Ft}[1]{{\widehat{#1}}}
\newcommand{\IFt}[1]{{\widecheck{#1}}}

\newcommand{\ot}{{\leftarrow}}

\newcommand{\weight}{{\mathrm{weight}}}
\newcommand{\poly}{{\mathrm{poly}}}
\newcommand{\Poly}{{\mathrm{Poly}}}

\newcommand{\diff}{\mathop{}\!\mathrm{d}}

\numberwithin{equation}{section}

\DeclareMathOperator{\Id}{Id}
\DeclareMathOperator{\ad}{ad}
\newcommand{\AD}{{\mathbf{ad}}}
\newcommand{\F}{{\mathcal{F}}}
\newcommand{\noset}{\varnothing}
\newcommand{\mbf}[1]{{\mathbf{#1}}}
\newcommand{\mcal}[1]{{\mathcal{#1}}}
\newcommand{\wtild}[1]{{\widetilde{#1}}}
\newcommand{\PathSet}{{\Pi}}
\newcommand{\ddt}{{\frac{d}{dt}}}

\DeclareMathOperator{\Tr}{Tr}

\DeclareMathOperator{\Op}{Op}
\DeclareMathOperator{\Rept}{Re}
\DeclareMathOperator{\Impt}{Im}
\DeclareMathOperator{\vol}{vol}
\DeclareMathOperator{\supp}{supp}
\DeclareMathOperator{\Variance}{Var}
\DeclareMathOperator{\argmax}{argmax}
\DeclareMathOperator{\trace}{tr}
\DeclareMathOperator{\Div}{div}

\DeclareMathOperator{\Arm}{arm}
\DeclareMathOperator{\Stick}{stick}

\newtheorem{theorem}{Theorem}
\numberwithin{theorem}{section}
\newtheorem{lemma}[theorem]{Lemma}
\newtheorem{definition}[theorem]{Definition}
\newtheorem{corollary}[theorem]{Corollary}
\newtheorem{proposition}[theorem]{Proposition}
\newtheorem{example}{Example}
\newtheorem{remark}[theorem]{Remark}

\newtheorem{hope}[theorem]{Potentially Provable Lemma}
\newtheorem{assumption}{Useful Assumption}
\hfuzz=20pt

\title{Correlation of fluctuations}
\date{\today}
\begin{document}
\maketitle

I want to look at
\[
\xp (\mcal{D}[R]-\theta) R.
\]
Write
\[
R = M - M (\lambda V - \lambda^2 \theta) R.
\]
So
\begin{align*}
\xp (\mcal{D}[R]-\theta)R
&= \xp (\mcal{D}[R]-\theta) M(\lambda V - \lambda^2\theta)R \\
&= \lambda^2 \sum_x\xp \mcal{D}[R\delta_zR] M\delta_z R
+ \xp (\mcal{D}[R]-\theta)M(\mcal{D}[R]-\theta)R.
\end{align*}
The crossing term is what I'm curious about.

Just for clarity, let's put an $M$ in the front, since that's usually where this comes up.
If we compute an entry, one thing we can do for example is bound like so:
\[
\lambda^2 (M (\mcal{D}[R]-\theta)R)_{00} \lsim
\lambda^2 \|M\|_{1\to 2} (\lambda) \|R\|_{1\to 2} \lsim \lambda
\]
(ok, taking $\eta\sim \lambda^2$).

I want to see that
\[
\lambda^4 \xp(M \sum_z \mcal{D}[R\delta_zR]M\delta_z R)_{00}
\]
is much smaller.  One way to go about this is to try to bound $\|X\|_{2\to 2}$ where
\[
X_{xy} := R_{xy}^2 M_{xy}.
\]
But this is pretty bad...
I guess what I want is that this is better than $\lambda^{-1}$, that would be the break-even with the previous bound.
Well, I could do
\[
\|X\|_{2\to 2} \leq \max_x \sum_y |R_{xy}|^2 |M_{xy}|
\leq (\sum_y |R_{xy}|^4)^{1/2} (\sum_y |M_{xy}|^2)^{1/2}.
\]
Well and I guess it's only the $\ell^4$ inside a kinetic ball, but that's still $O(1)$
(right?).

Could try instead for
\[
(\sum_y |R_{xy}|^3)^{2/3} (\sum_y |M_{xy}|^3)^{1/3}.
\]
And in high enough dimensions this is $O(1)$.


It's just weird that I am only seeing an improvement in very high dimensions.

\end{document}
